% Generated draft: requires mathematical and source review.
\documentclass[11pt]{article}
\usepackage[margin=25mm]{geometry}
\usepackage{fontspec}
\setmainfont{Latin Modern Roman}
\usepackage{amsmath,amssymb,mathtools,unicode-math}
\setmathfont{Latin Modern Math}
\usepackage{enumitem,needspace,xurl,hyperref}
\setlength{\parindent}{0pt}
\setlength{\parskip}{5pt}
\setlength{\emergencystretch}{3em}
\newcommand{\N}{\mathbb{N}}
\newcommand{\Z}{\mathbb{Z}}
\newcommand{\Q}{\mathbb{Q}}
\newcommand{\R}{\mathbb{R}}
\newcommand{\C}{\mathbb{C}}
\newcommand{\F}{\mathbb{F}}
\newcommand{\A}{\mathbb{A}}
\newcommand{\PP}{\mathbb{P}}
\newcommand{\E}{\mathbb{E}}
\newcommand{\eps}{\varepsilon}
\newcommand{\dd}{\,\mathrm{d}}
\begin{document}
\section*{UnsolvedMath ID: 800743}
\textbf{Source problem code:} catalogue entry
\subsection{Definitions and mathematical statement}
Let \(X\) be a Banach space over a scalar field \(\mathbb K\), and let \((e_n)_{n\geq 1}\) be a normalized Schauder basis: \(\|e_n\|=1\), and every \(x\in X\) has a unique norm-convergent expansion \(x=\sum_{n=1}^{\infty}a_ne_n\), with \(a_n\in\mathbb K\). The supplied record does not specify whether \(\mathbb K=\R\) or \(\mathbb K=\C\).

For the same scalar field, define
\[
c_0(\mathbb K)=\{(a_n)_{n\geq1}:a_n\in\mathbb K,\ a_n\to0\},\qquad \|(a_n)\|_\infty=\sup_n|a_n|,
\]
and, for \(1\leq p<\infty\),
\[
\ell_p(\mathbb K)=\left\{(a_n)_{n\geq1}:a_n\in\mathbb K,\ \sum_{n=1}^{\infty}|a_n|^p<\infty\right\},\qquad
\|(a_n)\|_p=\left(\sum_{n=1}^{\infty}|a_n|^p\right)^{1/p}.
\]
Two Banach spaces are isomorphic if there is a bounded linear bijection between them with bounded inverse.

\textbf{Source question (block homogeneous basis problem).} If \(X\) has a normalized block homogeneous Schauder basis, must
\[
X\cong c_0(\mathbb K)\quad\text{or}\quad \exists p\in[1,\infty)\; X\cong\ell_p(\mathbb K)?
\]
The hypothesis ``block homogeneous'' is retained from the source. Its definition is absent from the supplied record, so the question above is not yet a self-contained formulation.

To make explicit one possible block-subspace interpretation requiring source confirmation, a normalized block sequence of \((e_n)\) is a sequence \((x_k)_{k\geq1}\) of finitely supported vectors with \(\|x_k\|=1\) and
\[
\max\operatorname{supp}_{e}(x_k)<\min\operatorname{supp}_{e}(x_{k+1})\qquad(k\geq1),
\]
where \(\operatorname{supp}_{e}(\sum_n a_ne_n)=\{n:a_n\ne0\}\). If ``block homogeneous'' means
\[
\overline{\operatorname{span}}\{x_k:k\geq1\}\cong X
\quad\text{for every normalized block sequence }(x_k)_{k\geq1},
\]
then the question asks whether every \(X\) admitting a normalized Schauder basis with this property has the displayed conclusion. This expansion is conditional on verification of the source terminology; it is not asserted to be the definition intended by the unresolved reference [FR].
\subsection{Short English statement}
The supplied question asks whether a Banach space with a normalized block homogeneous Schauder basis must be isomorphic to \(c_0\) or to \(\ell_p\) for some \(1\leq p<\infty\). The meaning of block homogeneity and the intended scalar field require confirmation. The registry labels the problem open; that label is supplied metadata, not an independently verified status assessment.
\subsection{Sources}
\begin{itemize}
\item Local UnsolvedMath registry, permanent ID \(800743\), source code \texttt{catalogue entry}, title: \emph{Banach Spaces and Ramsey Theory --- Problem 6}. Supplied upstream revision: \texttt{d371ee917ac9647fbb06cd6fc529deec6f3dda82}.
\item Dodos and Lopez-Abad, \emph{Banach Spaces and Ramsey Theory: Some Open Problems} (2010), Problem 6, \url{https://arxiv.org/abs/1006.2668}. The supplied record attributes the problem to Dodos and Lopez-Abad, with proposed year 2010. Its provenance reports extraction from arXiv LaTeX and an access date of 2026-10-16; these are supplied provenance statements, not claims of independent source access.
\item The source statement calls this the ``Block homogeneous basis problem'' and says ``see [FR]''. The supplied record does not identify the bibliographic entry [FR], so that attribution cannot be expanded reliably.
\end{itemize}
\end{document}
